
@@M0@@
est
injectif, donc @@M1@@ l'est aussi.

6.5. Soit @@M2@@ un schéma noethérien de dimension @@M3@@, ayant un faisceau ample @@M4@@. On munit @@M5@@ de la filtration définie comme suit (cf. [1]): @@M6@@ appartient à @@M7@@ si et seulement si pour tout fermé @@M8@@ de @@M9@@, il existe un complexe @@M10@@, parfait, tel que:
@@M11@@
@@M12@@

Les @@M13@@ ainsi définis sont des sous-groupes additifs de @@M14@@. En effet, soient @@M15@@, et pour un fermé @@M16@@ donné dans @@M17@@, @@M18@@ et @@M19@@ les complexes correspondants; alors @@M20@@ est la classe du complexe @@M21@@, et
@@M22@@
comme la codimension d'une réunion finie de fermés est la borne inférieure des codimensions des fermés, on a bien @@M23@@.

\textbf{Proposition 6.6.} \uline{La filtration} 6.5 \uline{de} @@M24@@ \uline{est une filtration d'anneau} \uline{et} @@M25@@.

Soient @@M26@@, @@M27@@, et @@M28@@ un fermé de @@M29@@. Il existe un complexe @@M30@@ tel que @@M31@@, et @@M32@@; il existe un complexe @@M33@@ tel que @@M34@@, et :
@@M35@@

On a alors : @@M36@@, et :
@@M37@@
(EGA @@M38@@ 14.2.2.3).
